\noindent{\large\bfseries 10 月 22 日补充：sample3 原账本与 2024 年数据核对}
\par\noindent{\small 保留前文及 10 月 21 日补充。本次为既有策略的执行压力回放接入，不是样本外检验。}

\lead{这次组合与观察范围已经对上}
新包的日账覆盖 2019 年 1 月 31 日至 2024 年 8 月 30 日，共 1,355 行；其中 2024 年 1 月 2 日至 8 月 30 日为 162 行。与原研究 \texttt{run\_manifest} 中默认 $N=20$、sample3、\texttt{H20\_equal\_G1} 的日期、行数、末净值、买卖金额及费用汇总逐项一致。$N$ 是因子回看期，$H$ 是调仓周期（交易日），均为 20；重试与退出会改变实际持有期。实际方向为低因子 G1、等权。尚未独立重算逐股因子值、排名和分组，不能将汇总一致扩大为选股全流程验收。没有改成前文演示用的追涨前 15 只，也没有重新选参数。

2024 年属于参与过原参数选择的 sample3。后续即使补齐并完成回放，也只能按既有策略的执行压力回放解释，不能写成独立样本外收益或策略容量已获验证。

\lead{已完成金额账；原始股数与独立估值仍缺}
实际读取了 139,409 行成交载荷、131,863 行持仓事件和 1,339 行日级执行尝试汇总。逐证券金额汇总到日账及尝试表，最大差约 $4.44\times10^{-16}$；逐笔费用与买入 0.015\%、卖出 0.065\% 的公式一致。按来源归一化期初现金 1 连续递推，全期现金最大差约 $9.33\times10^{-15}$。这验证了金额和现金等式，没有把归一化账户直接解释为人民币或原始股数。

43,223 行卖出记录中有 36,687 行的 \texttt{units} 恰为零，说明该字段不能当卖出成交数量。随后提供的源码已确认它是执行后的复权持仓余额。原账户只有现金与复权持仓，按现金加持仓单位乘复权估值价计算净值；没有独立的原始股数、可卖量、负债及公司行动到账流水。资金尺度明确后，成交金额除以原始价可推导等价数量，但未必为合法整数股数，复权持仓也不能直接当真实可卖持仓，尤其不能再重复计现金分红。

\begin{table}[H]\centering\small
\begin{tabularx}{\linewidth}{p{31mm}Y}
\toprule 先后顺序选窗 & 已知日期及范围（仅为读取检查，未回放成交）\\\midrule
前两次新目标 & 1 月 11 日、2 月 8 日；后续持有观察到 3 月 14 日，共 40 行日账、5,196 行原成交。\\
前三次新目标 & 再含 3 月 15 日；后续观察到 4 月 15 日，共 60 行日账、7,077 行原成交。\\
窗口前账户 & 1 月 10 日日账报告持仓 1,315 只、现金为零，但证券字典、原始股数、可卖量及权益状态仍缺；还须取得 2023 年末完整账户以承接 1 月 2 日。\\
\bottomrule
\end{tabularx}
\caption{日期按实际 \texttt{new\_target} 顺序确定，没有根据收益挑窗；日期不等于已知指令时点。}
\end{table}

1 月 2 日至 10 日是上年目标的重试或股票池退出汇总，不能当本年新调仓。源码确认原计划是持续目标权重：新调仓覆盖旧目标与待执行方向；逐日按开盘净值、实际持仓重算需求，旧买卖方向不反转；退出股票池时目标归零且不重新分配。\texttt{orders.csv} 只有日级汇总，仍缺实际逐证券权重及更新状态，不能倒造固定股数母单。源码以各证券实际开盘价记成交，没有秒级指令时点；\texttt{pool\_exit\_count} 是当前退出状态计数，不是每天新增同样多的退出交易。

\clearpage
\lead{L2 已扩大覆盖，但相关数据仍未全部通过验收}
10 月 22 日通过现有授权只读接口，对 2024 年实际成交出现的 4,026 个证券、25 个成交日期，共 20,414 个证券日期，逐批读取 L2 与日线整日聚合。查询单并发、每次最多 200 行、30 秒，私有配置留在服务器。已有 MySQL 可读配置副本可用，无需读取或复制原秘密文件。

\begin{table}[H]\centering\small
\begin{tabularx}{\linewidth}{Y r}
\toprule 本次核查项 & 数量\\\midrule
实际成交涉及的证券日期 & 20,414\\
有 L2 记录的证券日期 & 20,410\\
上述 L2 原记录总行数 & 77,272,791\\
缺 L2（同一存托凭证，保留缺项） & 4\\
恰有一行日线的证券日期 & 20,414\\
含负成交额增量的证券日期 & 131\\
L2 日末累计额与日线金额差大于 1 元 & 1,809\\
L2 日末价格与日线收盘价差大于半分 & 2\\
\bottomrule
\end{tabularx}
\caption{差异分类可以重叠；这些是整日字段核对，不能代替盘中连续性、时点及可交易状态验收。}
\end{table}

\panel{06_sample3_data_checks}{2024 年实际成交证券日期的读取范围与原字段差异。图中没有正式策略成交完成率或收益。}

有 L2 的日期中，区间量额求和与日末累计字段一致，累计量与日线成交量差不超过 1 股；但这不证明中间每段完整。131 个较大的金额异常集中在 2 月 8 日和 6 月 17 日，并含负成交额增量；其余较小差异也保留原值待核。未自动把金额乘十、补区间量或以日终总量替代可提前使用的成交量。日末价差的两项亦未强行改写。

这次日线复权字段已可读，与此前部分检查的空字段状态不同；但还没有证明它与原账本共享批次一致。日线主要来自 L2 派生链，不是完全独立源。162 个日账日期与沪深交易日历一致，全部实际成交证券日期有涨跌停记录；历史 ST 视图返回 160 条正面记录，停牌视图为空，缺行均不自动判为正常。前三期候选证券读取到 13 条公司行动动作，仍须实际持仓与完整阶段覆盖才能计权益。此次范围不包括未知未成交目标、全部旧持仓和可能延续的需求，不能宣称完整账户行情覆盖已经通过。

\clearpage
\lead{跨市场数量规则已补入，正式输入仍须保持原策略语义}
新增可选 \texttt{cn\_stock\_2024} 模式，依据 2024 年适用的沪深历史规则；旧演示默认保留。规则限制新子单，不改给定母单数量；显式选择 \texttt{persistent} 时，同一合法子单部分成交后，在仍有效时保留余量身份。可卖数量和 T+1 仍由同一连续账户约束。

\begin{table}[H]\centering\small
\begin{tabularx}{\linewidth}{Y r r r}
\toprule 板块 & 买入最小／递增 & 限价最大 & 市价最大\\\midrule
沪深主板 & 100／100 & 1,000,000 & 1,000,000\\
创业板 & 100／100 & 300,000 & 150,000\\
科创板 & 200／1 & 100,000 & 50,000\\
已明确身份的科创存托凭证 & 200／1份 & 100,000份 & 50,000份\\
\bottomrule
\end{tabularx}
\caption{卖出零股按板块另行处理，不把买入最小量套在有效子单余量上。数量类别须显式指定。}
\end{table}

主板和创业板的持仓非整百余数须一次卖出，可与整百部分合并；科创板持仓不足 200 时须一次卖出全部。存托凭证须明确数量为份、价格为元／份，期初账户也须有单位声明。现有原券单适配器仍只接受原始股票股数，因此存托凭证尚未正式接通；其 L2 缺项及公司行动合同排除范围也必须另行解决。

依据为\href{https://www.sse.com.cn/lawandrules/sselawsrules2025/repeal/rules/c/c_20250612_10824490.shtml}{上交所 2023 年交易规则}第 3.3.8、3.3.9、6.1.7 条，\href{https://investor.szse.cn/institute/rules/t20230706_601604.html}{深交所 2023 年主板入市手册}，及\href{https://edu.sse.com.cn/tib/qa/c/4866268.shtml}{上交所科创板数量说明}。选择限价或市价数量上限，只定义申报量边界，不表示十档模型已经实现完整市价委托、保护限价或排队机制。

独立审查已修正科创板目标权重入口先按 100 股取整，以及存托凭证期初持仓无单位也可估值的问题。目前该模式只接固定母单；原目标权重尚缺来源明确的缩放、取整和净需求规则，因此拒绝沿用旧演示入口。全部 277 项测试通过，复现命令随图源及代码文档提供；合成恢复测试不冒充真实 sample3 恢复验收。

真实量额异常还揭示旧差分检查的遗漏：负额与累计额同步下降仍可通过。现已强制排除负额和非有限值，并新增显式 \texttt{price\_bounds} 模式，以当前快照已知日内高低价核查区间金额；正成交量缺界也不提供额度，不用未来日线。1 元仅为工程核查容差，不是行情精度保证。两个真实 10 行片段各保留 5 条正常来源，排除 1 条负额及随后 3 条十分之一金额，另 1 条因缺前快照不提供额度。原始量额未改，没有运行虚构订单。

\lead{继续正式回放仍缺哪些具体输入}
新提供的 \texttt{code\_snapshot} 已补齐读取实现并完成静态审阅；原读取器和导出脚本含自动校验流程，本项目没有调用它们。现已按明确数据结构独立适配：共享输入是原始开盘价、复权开盘价、复权因子、估值价、估值沿用标志五张宽表，证券顺序来自估值表列名；持仓事件是余额替换，不是成交股数增量。仍待新的 sample3 数据或同组合普通表，以实际恢复而非合成测试验收。

还需实际逐证券目标权重及其更新、退出和待执行状态，来源明确的指令时点，以及资金尺度、股级取整和复权账户与真实公司行动如何对应的约定。2023 年末及小窗口前完整账户、同期 100001 指数、相关状态与公司行动的完整覆盖须按真实持仓补齐。日级日期没有自行填为 09:30。

原账本含费，而主实验两侧均为零费。当前仅已验证固定已成交金额的退费现金路径，独立逐股估值与原始股数尚未恢复，不能输出原零费执行收益、L2 对照收益、两侧指数超额、收益差或资金容量。输入齐备后，先核对上述前两三期原账户，再给同一计划订单做 L2 对照；账户只初始化一次，之后承接实际持仓扩大到全期，再运行显式配置的资金阶梯。本次没有伪造母单、补成交或替换参数，也未做任何哈希校验。
